\documentclass[11pt,a4paper]{article}

\usepackage[margin=2.35cm]{geometry}
\usepackage{amsmath,amssymb,mathtools,bm}
\usepackage{booktabs,array}
\usepackage{enumitem}
\usepackage[numbers,sort&compress]{natbib}
\usepackage[colorlinks=true,linkcolor=blue,citecolor=blue,urlcolor=blue]{hyperref}
\setlength{\emergencystretch}{2em}
\sloppy

\newcommand{\dd}{\mathrm d}
\newcommand{\ii}{\mathrm i}
\newcommand{\Id}{\mathbb I}
\newcommand{\diag}{\operatorname{diag}}

\title{MadStree: tree-level dS integrals, \texorpdfstring{dlog}{dlog}-form differential equations and package usage}
\author{MadStree package documentation}
\date{August 3, 2026}

\begin{document}
\maketitle

\begin{abstract}
We present the formulas underlying the Wolfram Language package \texttt{MadStree} (v0.5), which
computes master integrals, iterative reduction data and \text{dlog}-form differential equations for
arbitrary massive/massless mixed de~Sitter (dS) tree integrals, and evaluate them numerically
through the Python/FLINT backend FlintNDE. The derivation follows the notation of
\citet{Chen:2023eic,Chen:2024jms}: vertex integrals are organized into Kronecker sums of
two-dimensional Hankel systems, the differential equations are written in the form
$\dd I=\dd\Omega\cdot I$, and the boundary conditions are set at the large-damping-energy limit by
nested blow-up. We first derive the formulas, then describe how to use the package.
\end{abstract}

\tableofcontents

\part{Basic formulas}

\section{Conventions and local function systems}
\label{sec:conventions}

We follow the conventions of \citet{Chen:2023eic}: the conformal time is $\tau<0$, $z=-k\tau$ for a
line of momentum $k$, and the two independent solutions of the Bessel equation are
\begin{equation}
  h(\nu,0;z)=z^{-\nu}H_{\nu}(z),
  \qquad
  h(\nu,1;z)=-\partial_z h(\nu,0;z),
  \label{eq:paper-h-definition}
\end{equation}
where $H_\nu$ is the Hankel function. MadStree fixes $\nu=|\nu|$, uses the single function name
$h$, and instead of \eqref{eq:paper-h-definition} defines
\begin{equation}
  h(\nu,0;z)=z^{\pm\nu}H_\nu(z),
  \label{eq:positive-order-h}
\end{equation}
where the plus sign is the default (\texttt{NuConvention->"Positive"}) and the minus sign
(\texttt{"Negative"}) is the paper convention.
The two bases differ by the factor $z^{2\nu}$. Because the Bessel equation contains only $\nu^2$,
all formulas of one convention are obtained from the other by the uniform replacement
\begin{equation}
  \left.\nu\right|_{2401}\longrightarrow-\nu,
  \label{eq:paper-to-positive-order}
\end{equation}
applied to every occurrence of $\nu$ (equations of motion, $M_1$, contact powers, recurrences and
the differential equations). Equivalently, the positive-prefactor basis is the paper's function with
negative order, since $H_{-\nu}\propto H_\nu$ up to a kinematics-independent phase.

For a vertex $v$ we write $k_{0,v}$ and $\nu_{0,v}$ for its damping energy and time power, and for a
line $e$ we write $k_e$ (or simply $k_i$) and $\nu_i$ for its momentum and Bessel order. The
propagator of a massive or massless full line is denoted by $G^{\sigma_e\sigma'_e}$ with
$\sigma_e=\pm$; we use $\sigma_e=+1,-1$ for $G^{++},G^{--}$. A theta-free exponential factor is
written as $e^{\ii\chi_bq_b\tau}$, with $\chi_b=\pm1$; its only role is to contribute to the signed
phase energy
\begin{equation}
  p_v=\sum_{b:\,v(b)=v}\chi_bq_b,
  \label{eq:vertex-phase-energy}
\end{equation}
so that the vertex time exponent is $e^{\ii p_v\tau}$. Pure vertex energies, massless cross/external
branches and auxiliary regulators are all such factors. The package's public time power is
$(-\tau_v)^{A_v}$, so that $\partial_{\tau_v}(-\tau_v)^{A_v}=-A_v(-\tau_v)^{A_v-1}$.

We now turn to the two-dimensional $h$-system, the second local building block, and to the
Kronecker-sum form of the single-vertex IBP.

\subsection{The two-dimensional $h$-system}

With $z=-k\tau$, the vector $\bm h=(h(\nu,0;z),h(\nu,1;z))^T$ obeys the first-order system
\begin{equation}
  \partial_\tau\bm h
  =\left[-\ii k\sigma_2-\frac{2\nu+1}{\tau}\frac{\Id-\sigma_3}{2}\right]\bm h,
  \label{eq:h-first-order-system}
\end{equation}
where $\sigma_2,\sigma_3$ are Pauli matrices. In the dimensionless variable $z$,
\begin{equation}
  \partial_z^2h(\nu,0;z)+\frac{2\nu+1}{z}\,\partial_zh(\nu,0;z)+h(\nu,0;z)=0,
  \label{eq:h-dimensionless-eom}
\end{equation}
with constant term $1$; the physical $k^2$ is recovered only through $\partial_\tau=-k\partial_z$.
For the default positive prefactor and $\nu=1/2$,
\begin{equation}
  z^{1/2}H_{1/2}^{(1,2)}(z)\propto e^{\pm\ii z},
  \qquad
  (2\nu+1)\big|_{2401}\xrightarrow{\nu\to-\nu}1-2\nu=0,
  \label{eq:massless-h-eigenstate}
\end{equation}
so the first-derivative state is a branch eigenstate, $h^{(1,2)}(1/2,1)=\pm\ii h^{(1,2)}(1/2,0)$.

\subsection{Vertex integrals and their Kronecker-sum IBP}

For a single vertex with $p$ massive/full legs, the integrals are
\begin{equation}
  V_v(\nu_{0,v},\bm a)
  =\int\dd\tau_v\,\tau_v^{\nu_{0,v}}e^{\ii p_v\tau_v}
  \prod_{i=1}^{p}h(\nu_i,a_i;-k_i\tau_v),
  \qquad a_i\in\{0,1\},
\end{equation}
and we collect them into the vector $\bm f_v(a_0)=(V_v(\nu_{0,v}+a_0,\bm a))_{\bm a}$. In the paper
\citet{Chen:2023eic}, the IBP is
\begin{equation}
  \left[M_1+\nu_{0,v}\Id\right]\bm f_v(-1)
  +\left[M_0+\ii p_v\Id\right]\bm f_v(0)=0,
  \label{eq:vertex-ibp}
\end{equation}
with the local matrices
\begin{equation}
  M_1=-\sum_{i=1}^{p}(2\nu_i+1)\frac{\Id-\sigma_3^{(i)}}{2},
  \qquad
  M_0=\sum_{i=1}^{p}(-\ii k_i\sigma_2^{(i)}),
  \label{eq:vertex-M1-M0}
\end{equation}
where $A^{(i)}$ acts on the two-dimensional factor of leg $i$ and trivially elsewhere. Both are
diagonalized by the constant matrix
\begin{equation}
  T=\frac1{\sqrt2}
  \begin{pmatrix}1&-\ii\\-\ii&1\end{pmatrix},
  \qquad
  T\sigma_2T^{-1}=\sigma_3,
  \label{eq:paper-T}
\end{equation}
in the sense that $T_pM_0T_p^{-1}$ with $T_p=T^{\otimes p}$ is diagonal and $M_1$ is already diagonal
in the binary basis. One lowering and one raising step are therefore
\begin{equation}
  A_{-}(\nu_0)=-M_1(\nu_0)^{-1}M_0,
  \qquad
  A_{+}(\nu_0)=-T_p^{-1}\widetilde M_0^{-1}T_p\,M_1(\nu_0+1),
  \label{eq:vertex-Apm}
\end{equation}
with $\widetilde M_0=T_pM_0T_p^{-1}$; only diagonal matrices are inverted.

\section{The full-graph building blocks and package conventions}
\label{sec:full-graph}

\subsection{From vertex factors to the full graph}

Each active massive endpoint is a two-dimensional $h$-slot, each uncontracted massless full edge a
shared two-dimensional slot (Sec.~\ref{sec:massless}), and each theta-free exponential a
one-dimensional factor contributing only to \eqref{eq:vertex-phase-energy}. The full-graph space of
a sector $s$ is the tensor product of these factors, with dimension $2^{|H_s|+|E_s|}$, where
$H_s$ are the active massive endpoint slots and $E_s$ the uncontracted massless edges. We call the
bookkeeping of slots and their vertex incidence the \emph{building-block registry}.

For a vertex $v$ the global regular matrices are the sums of the local atoms \eqref{eq:vertex-M1-M0}
embedded into the common space:
\begin{align}
  M_{1;v}
  =\nu_{0,v}\Id
  -\sum_{i\in H_v}(2\nu_i+1)\frac{\Id-\sigma_3^{(i)}}{2},
  \label{eq:global-M1}\\
  M_{0;v}
  =\ii p_v\Id
  +\sum_{i\in H_v}(-\ii k_i\sigma_2^{(i)})
  +\sum_{e\in E_v}(-\ii\epsilon_{ve}k_e\sigma_2^{(e)}),
  \label{eq:global-M0}
\end{align}
where $\epsilon_{ve}=+1,-1$ if $v$ is the first, second endpoint of the ordered edge $e$, and
$\sigma_2^{(e)}$ acts on the shared slot of $e$. The massless edge does not enter $M_1$, because the
default substitution \eqref{eq:paper-to-positive-order} gives $2\nu+1\to1-2\nu=0$ for
$\nu=1/2$. All $M_{0;v}$ commute pairwise, so a single constant
$U_s=\bigotimes T$ diagonalizes them, giving the energy letters
\begin{equation}
  \kappa_{v;\bm a,\bm r}
  =p_v
  +\sum_{i\in H_v}(2a_i-1)k_i
  +\sum_{e\in E_v}\epsilon_{ve}(2r_e-1)k_e,
  \label{eq:global-energy-letter}
\end{equation}
with $a_i,r_e\in\{0,1\}$ the binary state of the massive and shared slots. The same shared bit $r_e$
enters the letters of both endpoints with opposite signs. The common diagonalization is legitimate
only if, on every slot, the local operators of all vertices contain a single Pauli direction
($\sigma_2$ for massive and paper-$h$ shared slots, $\sigma_1$ for the package quotient of
Sec.~\ref{sec:massless}); otherwise the fast diagonal inversion below is disabled.

The remaining building block is the theta-carrying massless full edge. Before reduction it behaves
like a massive propagator, but its extra relations project the four endpoint states onto a single
shared two-state:

\subsection{Massless full edges: from four states to a shared two-state}
\label{sec:massless}

Before the extra relations are used, a theta-carrying massless full edge behaves like a massive one:
each endpoint carries a two-state $h$-factor, i.e.\ four states
$\bm F^{(4)}=(F_{00},F_{01},F_{10},F_{11})^T$. By \eqref{eq:massless-h-eigenstate}, in the
dimensionless paper basis
\begin{equation}
  F_{01}=-F_{10},
  \qquad
  F_{11}=F_{00},
  \label{eq:h-state-massless-relations}
\end{equation}
where $F_{ab}$ uses $h_1=-k^{-1}\partial_\tau h_0$. The quotient to the shared two-state
$\bm g=(F_{00},F_{10})^T$ is the constant projection
\begin{equation}
  \bm F^{(4)}=S_e\bm g,
  \qquad
  S_e=
  \begin{pmatrix}1&0\\0&-1\\0&1\\1&0\end{pmatrix},
  \qquad
  P_e=
  \begin{pmatrix}1&0&0&0\\0&0&1&0\end{pmatrix},
  \qquad
  P_eS_e=\Id_2,
  \label{eq:SP-massless}
\end{equation}
and any operator preserving the relation subspace reduces to $O_{\rm red}=P_eOS_e$. In particular,
\begin{equation}
  P_e(\sigma_2\otimes\Id_2)S_e=\sigma_2,
  \qquad
  P_e(\Id_2\otimes\sigma_2)S_e=-\sigma_2,
\end{equation}
so the two endpoints act on the \emph{same} shared slot with opposite signs. This is the only
regular reason why per-vertex factorization fails; the full-graph Kronecker-sum structure is
preserved. The four-state contact vector of Eq.~(3.65) of \citet{Chen:2023eic} is
$\bm r^{(4)}=(0,-1,1,0)^T=S_e|1\rangle_e$, so the contact selector in the shared basis is simply the
fixed column $|1\rangle_e=(0,1)^T$, with weights $w^{(u_e)}=+2\ii\sigma_e$,
$w^{(v_e)}=-2\ii\sigma_e$ in the paper basis.

The package uses an algebraic quotient that separates the endpoint-derivative phases from the
discrete states,
\begin{align}
  \mathsf m_{e,0}
  &=\theta(\tau_{u_e}-\tau_{v_e})
    e^{-\ii\sigma_ek_e(\tau_{u_e}-\tau_{v_e})}
    +\theta(-(\tau_{u_e}-\tau_{v_e}))
    e^{+\ii\sigma_ek_e(\tau_{u_e}-\tau_{v_e})},\\
  \mathsf m_{e,1}
  &=-\theta(\tau_{u_e}-\tau_{v_e})
    e^{-\ii\sigma_ek_e(\tau_{u_e}-\tau_{v_e})}
    +\theta(-(\tau_{u_e}-\tau_{v_e}))
    e^{+\ii\sigma_ek_e(\tau_{u_e}-\tau_{v_e})}.
\end{align}
Their relations are $F^J_{01}=-F^J_{10}$, $F^J_{11}=-F^J_{00}$. The two shared bases are related by
the constant $D_{\sigma_e}=\diag(1,\ii\sigma_e)$, under which the endpoint generators become
$Q^J_{u_e}=+\ii\sigma_ek_e\sigma_1$ and $Q^J_{v_e}=-\ii\sigma_ek_e\sigma_1$, and the contact weights
become the exact numbers $w^J_{(u_e)}=-2$, $w^J_{(v_e)}=+2$. The package diagonalizer for these
$\sigma_1$ slots is the self-inverse Hadamard matrix
\begin{equation}
  H=\frac1{\sqrt2}\begin{pmatrix}1&1\\1&-1\end{pmatrix},
  \qquad
  H\sigma_1H^{-1}=\sigma_3,
\end{equation}
so both bases produce the same energy letters \eqref{eq:global-energy-letter}. Every regular and
contact matrix of the paper basis transforms to the package quotient by the constant similarity
$O^J=B_sO^hB_s^{-1}$ (and $C^J_{st}=B_sC^h_{st}B_t^{-1}$ for contacts), with
$B_s=\bigotimes D_{\sigma_e}$ on the massless slots.

\section{Iterative reduction and the dlog-form differential equation}
\label{sec:reduction}

The remaining formulas are the two outputs of the package: the iterative reduction and the
dlog-form differential equation.

\subsection{Iterative reduction}

The package's native integrals are
\begin{equation}
  J_s(\bm n;\bm a)=\texttt{MSIntegral[}s,\bm n,\bm a\texttt{]},
  \qquad \bm n\in\mathbb Z^{c_s},\ \bm a\in\{0,1\}^{N_s},
  \label{eq:MSIntegral-domain}
\end{equation}
where $c_s$ is the number of time components and $N_s$ the number of two-dimensional slots of the
sector. The sector $s$ is identified by the fixed-length bit string of the root propagators
(contracted $=0$, uncontracted $=1$; the top sector is all ones). The binary state order is
$\bm a(m)=\operatorname{IntegerDigits}(m,2,N_s)$; deleting a contracted endpoint bit and inserting
its complement at the other endpoint is done explicitly in the contact maps.

Writing the theta remaining term as a linear combination of lower-sector vectors
$R_{v,s}(\bm n)=\sum_{e\ni v}C_{s,s/e}^{(v,e)}\bm f_{s/e}(\bm n_e)$, with $\bm n_e$ the child time
shifts obtained by merging the shifts of the two endpoints, the one-step IBP at vertex $v$ is
\begin{equation}
  M_{1;v}(A_v+n_v)\bm f_s(\bm n-\bm e_v)
  +M_{0;v}\bm f_s(\bm n)
  +R_{v,s}(\bm n)=0,
  \label{eq:global-one-step-ibp}
\end{equation}
giving the lowering and raising steps
\begin{align}
  \bm f_s(\bm n-\bm e_v)
  &=A_{-,v}(\bm n)\bm f_s(\bm n)
  -M_{1;v}^{-1}R_{v,s}(\bm n),
  &A_{-,v}(\bm n)&=-M_{1;v}(A_v+n_v)^{-1}M_{0;v},
  \label{eq:global-Aminus}\\
  \bm f_s(\bm n+\bm e_v)
  &=A_{+,v}(\bm n)\bm f_s(\bm n)
  -U_s^{-1}\widetilde M_{0;v}^{-1}U_s
  R_{v,s}(\bm n+\bm e_v),
  &A_{+,v}(\bm n)&=-U_s^{-1}\widetilde M_{0;v}^{-1}U_s
  M_{1;v}(A_v+n_v+1).
  \label{eq:global-Aplus}
\end{align}
Lowering has a pole exactly when the $M_1$ eigenvalue
$\lambda_{1,v}(\bm a)=A_v+n_v-\sum_i(2\nu_i+1)a_i$ vanishes, and raising has a pole exactly on the
energy-letter surface $\kappa_{v;\bm a,\bm r}=0$; the zero matrix on the non-inverted direction
never produces a pole. Pure massless vertices have $M_1=\alpha_v\Id$ with
$\alpha_v=A_v+n_v$, and $\alpha_v=0$ corresponds to the unregularized divergence of the defining
$\Gamma(\alpha_v)$ integral; such cases are first regularized analytically
($\alpha_v\to\alpha_v+\delta$), reduced in the generic system, and expanded at $\delta\to0$.

If a bundle of full lines shares one theta, one derivative produces one delta, and the product
difference is a sum over odd nonempty subsets $S$ of the bundle with coefficient $2^{1-|S|}$; each
such subset is a simultaneous contact event. The base power of the merged component $C$ is
\begin{equation}
  A_C=\sum_{v\in C}A_v+\sum_{e\in E(C)}r_e+|C|-1,
  \label{eq:simultaneous-contact-base-power}
\end{equation}
with $r_e$ the raw contact power of the contracted edge. \texttt{MSReduce} iterates
\eqref{eq:global-Aminus}--\eqref{eq:global-Aplus} along the contact DAG until all shifts vanish,
and reports direction-sensitive singular layers.

The same building blocks also produce the differential equations. Differentiating the master
vector with respect to the kinematic variables and reducing the result with the recurrence above
yields the dlog-form potential in closed form:

\subsection{The \texorpdfstring{dlog}{dlog}-form differential equation}
\label{sec:dlog}

Following \citet{Chen:2024jms}, the master integrals satisfy
\begin{equation}
  \dd I=\dd\Omega\cdot I,
  \qquad
  \dd\Omega=\sum_i\Omega_{k_i}\dd k_i,
\end{equation}
with the diagonal matrices
\begin{equation}
  \widetilde\Omega_0=\diag\left[-\ii\log\kappa_{v;\bm a,\bm r}\right],
  \qquad
  \Omega_{\rm ex}=\diag\left[-\sum_i a_i(2\nu_i+1)\log k_i\right],
\end{equation}
and the fixed-sector potential
\begin{equation}
  \Omega_{ss}^{\rm reg}
  =\Omega_{\rm ex}
  -\ii\sum_{v\in V_s}
  U_s^{-1}\widetilde\Omega_0 U_s
  M_{1;v}(A_v+1),
  \label{eq:global-diagonal-omega}
\end{equation}
which is the full-graph version of Eq.~(2.12) of \citet{Chen:2024jms}. Its letters are exactly the
energy letters \eqref{eq:global-energy-letter}; no linear solution of IBP relations is needed.

The off-diagonal blocks use the remaining term $R^{(1)}$ of the raising relation,
\begin{equation}
  M_{0;v}\bm f_s(\bm e_v)=M_{1;v}(A_v+1)\bm f_s(0)-R_{v,s}(\bm e_v),
\end{equation}
and generalize Eq.~(3.68) of \citet{Chen:2023eic} to
\begin{equation}
  \Omega_{st}^{\rm contact}
  =+\ii\sum_{v\in V_s}\chi_v(E_{st})
  U_s^{-1}\widetilde\Omega_0 U_s
  C_{st}^{(v,1)},
  \qquad t\prec s,
  \label{eq:global-offdiagonal-omega}
\end{equation}
where $\chi_v(E_{st})$ is the sign of the vertex's pure energy block for pure-massive events, and
$\chi_v(E_{st})=(-1)^{N_0(E_{st})}$ when the event contains massless edges, $N_0$ being their
number. For a contracted massless edge with shared bit position $p_e$, the intermediate-state sum
collapses and the top-to-sub block is
\begin{equation}
  \left(\Omega_{s,s/e}\right)_{\bm a,\widehat{\bm c}}
  =2\ii\Big[
  (K_{u_e,s})_{\bm a,\bm c'}
  -(K_{v_e,s})_{\bm a,\bm c'}
  \Big],
  \qquad
  K_{v,s}=U_s^{-1}\widetilde\Omega_0 U_s,
  \label{eq:massless-binary-top-to-sub}
\end{equation}
up to a prefactor collecting the absorbed sector normalization and pinch factors, where $\bm c'$ is
the child state with the shared bit set to $1$. If $R^{(1)}$ lands on a nonzero child shift, the
program reduces each child state by the same recurrence and only writes the column into the
connection when all columns close; otherwise it returns \texttt{contactShiftReductionFailed}.
Ordering sectors by the number of remaining full edges, the complete connection is block
triangular, with contacts pointing only from a parent to a strictly lower sector.

The physical endpoint derivatives satisfy $\mathcal F_{01}=-\mathcal F_{10}$,
$\mathcal F_{11}=k^2\mathcal F_{00}$; if the final basis uses them instead of the dimensionless
$h$-states, the potential must be transformed by the gauge
\begin{equation}
  \mathcal A_{\rm phys}=\dd N\,N^{-1}+N\mathcal A_hN^{-1},
  \qquad
  N=\diag(1,-k),
\end{equation}
which contains $\dd\log k$ on the second entry.

It remains to fix the conventions that connect these formulas to the package interface.

\subsection{Package conventions}
\label{sec:package}

The package's public integrals are normalized objects $J_s=\mathcal N_sI_s$ with time powers
$(-\tau)^A$. In the normalized basis the sector diagonal block acquires $\dd\log\mathcal N_s\Id_s$
and the contact blocks are multiplied by $\mathcal N_s/\mathcal N_t$. Only when this ratio is a
kinematics-independent constant is \eqref{eq:global-offdiagonal-omega} directly a
constant-residue dlog. \texttt{NuConvention} is fixed at initialization. The
\texttt{MSInitVertexFamily} input uses the compact \texttt{ki/nui} form:
\begin{verbatim}
ctx = MSInitVertexFamily[<|
  "ki" -> {k0,k1,...}, "nui" -> {A0,nu1,...},
  "hankelBranches" -> {1,2,...}
|>];
\end{verbatim}
where \texttt{k0},\texttt{A0} are the vertex phase energy and base time power, and the remaining
entries give the $h$-block momenta and orders.

\section{Boundary conditions by nested blow-up}
\label{sec:boundary}

The last ingredient for numerical evaluation is the boundary data at the large-damping-energy
limit. With $\tau_v=-t_v<0$ and the damping energy $K_v=\ii p_v>0$, the vertex exponent is
$e^{\ii p_v\tau_v}=e^{-K_vt_v}$, so $K_v\to\infty$ localizes the integral to $\tau_v\to0^-$.
Choosing a strict total order of the damping energies,
\begin{equation}
  K_{\sigma(1)}\gg K_{\sigma(2)}\gg\cdots\gg K_{\sigma(V)}\gg1,
\end{equation}
defines the nested coordinates
\begin{equation}
  x_j=\frac{K_{\sigma(j+1)}}{K_{\sigma(j)}}\quad(1\leq j<V),
  \qquad
  x_V=\frac1{K_{\sigma(V)}},
\end{equation}
under which every theta becomes $0$ or $1$ ($\theta(\tau_u-\tau_v)\to1$ iff $K_u\gg K_v$). Every
$K$-dependent denominator of the connection has the affine form
$D_{\alpha,s}=Q_{\alpha,s}(k)+\sum_v c_{\alpha,s;v}K_v$ with kinematics-independent coefficients;
substituting $1/K_{\sigma(j)}=\prod_{r\ge j}x_r$ factors it into a coordinate monomial times a
nonzero unit, which removes the degenerate $K_i\to\infty$ singularities simultaneously. The program
verifies this explicitly: \texttt{MSBoundaryChartCertificate} aggregates all letters and
normalization divisors, extracts the coordinate monomial powers, requires the strict transforms to
be nonzero at $x=0$, and checks the normal-crossing Jacobian condition.
\texttt{RankOrder -> All} enumerates every strict chart, and \texttt{MSBoundaryData} refuses to
continue unless the chosen chart is certified.

The boundary data at $K_v\to\infty$ are the Frobenius leading branches $\{a,b,C\}$ of the master
integrals, obtained per sector from the indicial system
\begin{equation}
  (\lambda_{s,\bm a}\Id-R_0)v_{s,\bm a}=0
\end{equation}
of the pulled-back connection, with physical weights built from the component base powers,
normalizations and endpoint coefficients (the two-vertex $G^{++}$ example of
\citet{Chen:2024jms}, Sec.~4, is reproduced as a special case). FlintNDE then transports each exact
branch from the regular-singular point to a finite match point, sums the branches with their
weights, and transports along the affine pullback of the dlog connection to the user's ordinary
point. Save points declared by \texttt{FlintNDESavePoints -> \{\{coord,"save"\},...\}} are written
immediately; regular-singular starts store reusable $\{a,b,C\}$ data. Uncertified boundary charts,
unclosed shift reductions and backend capabilities beyond exact regular-singular transport all
fail closed.

\part{Using the package}

\section{Overview, requirements and loading}
\label{sec:overview}

MadStree is a topology-driven formula producer that implements the formulas of Part I. It accepts
\begin{itemize}
  \item an ordered dS tree topology with massive and/or massless full lines,
  \item a pure time-only incidence graph (cycles allowed, but no loop-momentum integration), or
  \item a single-vertex function family without a fake graph,
\end{itemize}
and produces, without first generating a large system of integration-by-parts (IBP) identities,
\begin{itemize}
  \item the contact-reachable sectors and their ordered master integrals
        $J_s(\bm n;\bm a)$ of \eqref{eq:MSIntegral-domain},
  \item single-step and complete iterative reduction data (Sec.~\ref{sec:reduction}),
  \item the block-triangular dlog-form potential $\Omega$
        (\eqref{eq:global-diagonal-omega}--\eqref{eq:global-offdiagonal-omega}),
  \item Frobenius boundary data at the $K_v\to\infty$ limit (Sec.~\ref{sec:boundary}),
  \item numerical master-integral values at user-specified ordinary points, through FlintNDE.
\end{itemize}
MadStree does not run loop-momentum integrations, Kira, or external reduction; anything that cannot
be certified (boundary charts, shift reductions, backend singular-point capabilities) fails closed
instead of degrading silently.

Before using the package we collect the system requirements and the loading commands.

\subsection{Requirements and loading}
\label{sec:loading}

\subsection{System requirements}
\begin{itemize}
  \item A Wolfram Language kernel: Mathematica or the free Wolfram Engine (the package is a
        \texttt{.wl} package; no external Wolfram libraries are needed).
  \item Python 3.10 or newer with \texttt{python-flint>=0.6} installed, for the numerical backend.
        MadStree invokes Python through the \texttt{PythonExecutable} option (default command
        \texttt{"python"}).
\end{itemize}

\subsection{Loading the package}

The version root is \texttt{.../package-MadStree/versions/MadStree-v0.5}. Add it to
\texttt{\$Path} and load with \texttt{Needs}:
\begin{verbatim}
madStreeVersionDir = ".../package-MadStree/versions/MadStree-v0.5";
AppendTo[$Path, madStreeVersionDir];
Needs["MadStree`"];
\end{verbatim}
Equivalently, load the kernel file directly,
\begin{verbatim}
$CharacterEncoding = "UTF-8";
Get[".../package-MadStree/versions/MadStree-v0.5/Kernel/MadStree.wl"];
\end{verbatim}
or use the official entry \texttt{Get[".../package-MadStree/load\_current.wl"]}, which loads the
version recorded in \texttt{VERSION\_INDEX.md} (currently v0.5). After loading,
\texttt{MSFlintNDEConfiguration[]} returns an association whose \texttt{availableQ} must be
\texttt{True}; it confirms that the embedded FlintNDE backend under \texttt{Vendor/FlintNDE} has
been found.

\section{Input and core computation}
\label{sec:inputs}

\subsection{Tree topology: \texttt{MSInitTree}}

A tree is given by an association of vertices and lines,
\begin{verbatim}
spec = <|
  "name" -> "masslessFullEdge",
  "vertices" -> {
    <|"id" -> v1, "energy" -> k1, "timePower" -> a1|>,
    <|"id" -> v2, "energy" -> k2, "timePower" -> a2|>
  },
  "lines" -> {
    <|"id" -> e1, "type" -> "masslessFull", "endpoints" -> {v1, v2},
      "momentum" -> q, "skType" -> "++", "nu" -> 1/2|>
  }
|>;
context = MSInitTree[spec];
\end{verbatim}
Vertex fields:
\begin{center}
\begin{tabular}{@{}lp{5.5cm}p{3.8cm}@{}}
\toprule
field & meaning & note\\
\midrule
\texttt{id} & vertex symbol & required\\
\texttt{energy} & damping energy variable $k_{0,v}$ (e.g.\ \texttt{k1}) & required\\
\texttt{timePower} & time-power variable $A_v$ (e.g.\ \texttt{a1}) & required\\
\texttt{phaseSign} & vertex exponent sign & optional; default $+1$\\
\bottomrule
\end{tabular}
\end{center}
Line fields:
\begin{center}
\begin{tabular}{@{}p{4.4cm}p{5.3cm}p{3.8cm}@{}}
\toprule
field & meaning & note\\
\midrule
\texttt{id} & line symbol & required\\
\texttt{type} & \texttt{"massiveFull"},\ \texttt{"masslessFull"},\\
                \texttt{"massiveExternal"},\ \texttt{"massiveCross"},\\
                \texttt{"masslessCross"} & required\\
\texttt{endpoints} & ordered pair of vertex ids & required\\
\texttt{momentum} & momentum variable $k_e$ (e.g.\ \texttt{q}) & required\\
\texttt{skType} & \texttt{"++"} or \texttt{"--"} & required for full lines\\
\texttt{nu} & Bessel/Hankel order $\nu=|\nu|$ & required for massive/full lines\\
\texttt{masslessRepresentation} & \texttt{"Quotient"} (default)\\ & or \texttt{"RedundantH"} &
                masslessFull only\\
\bottomrule
\end{tabular}
\end{center}
The initialization option \texttt{NuConvention -> "Positive"} (default) or \texttt{"Negative"}
chooses the prefactor sign of $h(\nu,0;z)=z^{\pm\nu}H_\nu(z)$ introduced in
Sec.~\ref{sec:conventions}; it is fixed for the whole context and cannot be changed per call.

\subsection{Time-only graphs: \texttt{MSInitTimeGraph}}

For a pure time-only graph the vertex and line fields are the same, but the graph may contain
cycles:
\begin{verbatim}
cycleContext = MSInitTimeGraph[<|
  "vertices" -> {<|"id" -> t1, "energy" -> k1, "timePower" -> a1|>, ...},
  "lines" -> {<|"id" -> l12, "type" -> "masslessFull", "endpoints" -> {t1, t2},
               "momentum" -> q12, "skType" -> "++", "nu" -> 1/2|>, ...}
|>];
\end{verbatim}

\subsection{Single-vertex families: \texttt{MSInitVertexFamily}}

No fake topology is needed for a single vertex. The compact input follows the reference
\texttt{ki/nui} convention,
\begin{verbatim}
vertexContext = MSInitVertexFamily[<|
  "ki" -> {k0, k1, k2},
  "nui" -> {a0, nu1, nu2},
  "hankelBranches" -> {1, 2}
|>, NuConvention -> "Positive"];
\end{verbatim}
where \texttt{First[ki]} is the pure vertex phase energy and \texttt{First[nui]} the base power of
$(-\tau)$; the remaining entries give, position by position, the $h$-block momenta and
$\nu=|\nu|$. The explicit alternatives \texttt{k0/timePower/hBlocks/exponentialBlocks} are
accepted.

Once the context is initialized, the following functions expose the formula products of Part I.

\subsection{Core computation}
\label{sec:core}

All functions take the initialized \texttt{context}.

\begin{center}
\begin{tabular}{@{}p{7.5cm}p{7.6cm}@{}}
\toprule
function & purpose\\
\midrule
\texttt{MSSectors[context]} & contact-reachable sectors in deterministic order (keys, contracted
lines, master counts)\\
\texttt{MSMasterIntegrals[context]} & ordered master integrals
\texttt{MSIntegral[sectorKey, timeShifts, stateBits]}\\
\texttt{MSFormulaMatrices[context, key]} & per-component $M_1$, $M_0$, $U$, $U^{-1}$, energy
letters $\kappa$\\
\texttt{MSDLogDE[context]} & block-triangular dlog potential $\Omega$, letters, certification
status\\
\texttt{MSReduce[expr, context]} & iterative reduction to the master basis\\
\texttt{MSContactMaps[context, key]} & parent-to-subsector contact maps\\
\texttt{MSWriteFormulaArtifacts[context]} & write formula products beside the caller\\
\bottomrule
\end{tabular}
\end{center}

\begin{description}
  \item[\texttt{MSSectors[context]}] returns the sector list. Each sector key is the fixed-length
        bit string of Sec.~\ref{sec:reduction}: \texttt{"0"} contracted, \texttt{"1"}
        uncontracted, in root propagator order; the top sector is all ones. Keys are strings, not
        integers; leading zeros are part of the identity.
  \item[\texttt{MSMasterIntegrals[context]}] returns the ordered master integrals
        \texttt{MSIntegral[sectorKey, timeShifts, stateBits]}, i.e.\ the $J_s(\bm n;\bm a)$ of
        \eqref{eq:MSIntegral-domain}. Their order is the unique authority for all matrices,
        boundary vectors and numerical results.
  \item[\texttt{MSFormulaMatrices[context, key]}] returns, per vertex component, the matrices
        $M_1$, $M_0$, the diagonalizer $U$ of \eqref{eq:global-M1}--\eqref{eq:global-M0}, its
        inverse, and the energy letters $\kappa$ of \eqref{eq:global-energy-letter}.
  \item[\texttt{MSDLogDE[context]}] returns \texttt{omegaPotential}, the matrix $\Omega$ of
        Sec.~\ref{sec:dlog} whose derivative along a path is the DE matrix, together with the
        letters, \texttt{letterMatrices}, \texttt{dlogResidual}, \texttt{constantResidueQ} and the
        certification \texttt{dlogStatus}. Only \texttt{"certifiedByFormulaChecks"} is a closed
        dlog and may be consumed numerically.
  \item[\texttt{MSReduce[expr, context, MasterBasis -> Automatic]}] reduces finite linear
        combinations of legal \texttt{MSIntegral}s by
        \eqref{eq:global-Aminus}--\eqref{eq:global-Aplus} along the contact DAG. Results include
        \texttt{result}, \texttt{masterBasis}, \texttt{coefficientVector}, \texttt{masterRules},
        \texttt{nonMasterResidual}, \texttt{remainingShiftedIntegrals} and
        \texttt{singularLayers}; a full reduction requires a zero residual and an empty shifted
        list.
\end{description}

\section{Boundary data and numerical evaluation}
\label{sec:numerics}

With an initialized context, numerical evaluation proceeds in two steps: boundary generation and
transport.

\subsection{Generating the boundary}

\begin{verbatim}
targetRules = {k1 -> -9 I, k2 -> -3 I, q -> 1, a1 -> 1, a2 -> 1};

certificate = MSBoundaryChartCertificate[context, targetRules, RankOrder -> All];

boundary = MSBoundaryData[context, targetRules,
  BoundaryScale -> 4, WorkingPrecision -> 40];
\end{verbatim}
\texttt{MSBoundaryChartCertificate} verifies the nested blow-up chart of Sec.~\ref{sec:boundary}
(normal crossing, theta fixing, quotient checks) for the chosen strict energy rank;
\texttt{RankOrder -> All} checks every strict chart. \texttt{MSBoundaryData} consumes the
certificate and returns the Frobenius leading branches $\{a,b,C\}$ at the $K_v\to\infty$ boundary.
If the chart is not certified, it returns \texttt{BoundaryChartNotCertified} and does not proceed
to numerics.

\subsection{Numerical transport}

\begin{verbatim}
targetValue = MSEvaluateTree[context, targetRules,
  BoundaryScale -> 4,
  WorkingPrecision -> 40,
  TransportOrder -> 80,
  ReferenceTransportOrder -> 104,
  TargetRelativeError -> "1e-20"];

targetValue["values"]
\end{verbatim}
For a single-vertex family use \texttt{MSEvaluateVertexFamily[vertexContext, targetRules]}.
Principal options:
\begin{center}
\begin{tabular}{@{}p{5.2cm}p{9.6cm}@{}}
\toprule
option & meaning\\
\midrule
\texttt{BoundaryScale} & distance of the finite match point from the boundary ($>1$)\\
\texttt{BoundarySeriesOrder} & truncation order of the boundary series\\
\texttt{RankOrder} & vertex order by damping energy $K_v$ (default: from the target point)\\
\texttt{WorkingPrecision} & working precision digits for the backend\\
\texttt{TransportOrder} / \texttt{ReferenceTransportOrder} & local-series orders of the
primary/reference chains\\
\texttt{TargetRelativeError} & target relative error between the two chains\\
\texttt{PythonExecutable} & Python command for the FlintNDE adapter (default \texttt{"python"})\\
\texttt{FlintNDESavePoints} & unnamed save points \texttt{\{\{coord,"save"\},...\}}\\
\texttt{MSRuntimeDirectory} & caller directory for temporary files (default: script directory)\\
\bottomrule
\end{tabular}
\end{center}

\subsection{Reading the result}

\texttt{targetValue} is an association with \texttt{status} (must be \texttt{"computed"}),
\texttt{values} (the master vector at the target, in master order), \texttt{masters},
\texttt{masterDigest}, \texttt{boundary}, \texttt{pathRules}, and \texttt{flintNDE}. For the
generic singular route, \texttt{flintNDE} has schema
\texttt{"madstree\_flintnde\_singular\_then\_ordinary\_v1"} with \texttt{singularLaunch} and
\texttt{ordinaryTransport}; the primary/reference relative difference is then inside
\texttt{ordinaryTransport}. A route-robust access is
\begin{verbatim}
relativeDifferenceInf = Lookup[
  targetValue["flintNDE"], "relativeDifferenceInf",
  Lookup[targetValue["flintNDE", "ordinaryTransport"],
         "relativeDifferenceInf", Missing["NotAvailable"]]];
targetValue["flintNDE", "targetRelativeErrorMet"]
\end{verbatim}

\section{Outputs, examples, validation and common pitfalls}
\label{sec:outputs}

\subsection{Output organization}

\begin{itemize}
  \item Transient JSON and Python caches: \texttt{results\_temp/\allowbreak flintnde\_transport/}
        beside the calling script; deleted after success, kept on failure for diagnostics.
  \item Save points: \texttt{results/\allowbreak flintnde\_save\_points/\allowbreak run-UUID/};
        each \texttt{\{coordinate,"save"\}} point is written immediately
        (\texttt{flintnde\_save\_001.json}, \dots), and the aggregate
        \texttt{madstree\_flintnde\_save\_points\allowbreak .json} only after full success.
  \item Formula artifacts: \texttt{MSWriteFormulaArtifacts[context]} writes \texttt{masters.wl},
        \texttt{recurrence\_metadata.wl}, \texttt{dlog\_de.wl} and \texttt{manifest.wl} to
        \texttt{results/\allowbreak madstree\_formula/\allowbreak run-UUID/}.
\end{itemize}
The package source directory never receives run artifacts.

The distribution ships three examples that exercise the complete workflow.

\subsection{Examples}
\label{sec:examples}

The v0.5 distribution ships three runnable examples under
\texttt{versions/MadStree-v0.5/Examples/}:
\begin{itemize}
  \item \texttt{01\_massless\_full\_edge.wl}: the massless quotient of Sec.~\ref{sec:massless},
        masters, recurrence, dlog and the automatic boundary/numerical entry;
  \item \texttt{02\_vertex\_family\_reduction.wl}: single-vertex family input, local tensor
        inverses, and finite-combination reduction;
  \item \texttt{03\_time\_only\_cycle\_chart.wl}: time-only cycle input, common-theta contacts, and
        all strict-rank chart certificates.
\end{itemize}
Run them in a Mathematica front end section by section, or with
\begin{verbatim}
wolframscript -file .../Examples/01_massless_full_edge.wl
\end{verbatim}

The following validation status and fail-closed guarantees define the supported domain of this
release.

\subsection{Validation status and limitations}
\label{sec:validation}

\begin{itemize}
  \item v0.5 passed all eleven development tests and the three examples (exit code $0$).
  \item Independent validations T1--T6 passed with counts $24/24$, $12/12$, $18/18$, $15/15$,
        $17/17$, $16/16$; T3 (the two-vertex $G^{++}$ system of Sec.~\ref{sec:boundary})
        reproduces the independent 2411 route with exactly zero connection residuals.
  \item The embedded FlintNDE test suite passes $76/76$.
\end{itemize}
Limitations and fail-closed guarantees:
\begin{itemize}
  \item No loop-momentum integration, no sampled IBP, no Kira or external reduction.
  \item \texttt{dlogStatus} must be \texttt{"certifiedByFormulaChecks"}; otherwise the DE is not
        consumed.
  \item Uncertified boundary charts, unclosed shift reductions, and singular points beyond the
        backend's certified capabilities all fail closed; ordinary Taylor transport is never used
        across singular points.
  \item The production boundary never uses defining integrals; finite-point defining integrals
        exist only inside independent validation programs.
\end{itemize}

Finally, we collect the most common mistakes.

\subsection{Common pitfalls}
\label{sec:pitfalls}

\begin{enumerate}
  \item \textbf{Convention}: \texttt{NuConvention} is chosen at initialization only
        (\texttt{"Positive"} by default). To compare with the formulas of
        \citet{Chen:2024jms}, initialize with \texttt{NuConvention -> "Negative"}, which selects
        the minus sign in \eqref{eq:positive-order-h}.
  \item \textbf{Sector keys}: fixed-length strings; leading zeros and the string type are part of
        the identity (Sec.~\ref{sec:reduction}).
  \item \textbf{Save points}: only \texttt{\{\{coordinate,"save"\},...\}} is accepted; no names,
        and coordinates must be exact Gaussian rationals.
  \item \textbf{Backend environment}: the Python invoked by \texttt{PythonExecutable} must have
        \texttt{python-flint} installed and must be able to import the embedded \texttt{flintnde};
        otherwise transport returns \texttt{FlintNDEProcessFailed} (inspect the failure's
        \texttt{backend} field or the retained \texttt{results\_temp} JSON).
  \item \textbf{Result access}: for the singular route, \texttt{relativeDifferenceInf} lives inside
        \texttt{ordinaryTransport}; use the fallback lookup of Sec.~\ref{sec:numerics}.
  \item \textbf{Encoding}: if loading fails with syntax errors, set
        \texttt{\$CharacterEncoding = "UTF-8"} before loading.
\end{enumerate}

We close with a brief summary of the package's data flow. MadStree implements the formulas above as
a single data flow: the building-block registry directly produces the Kronecker-sum matrices
\eqref{eq:global-M1}--\eqref{eq:global-M0}, the massless quotient \eqref{eq:SP-massless}, the
iterative recurrence \eqref{eq:global-Aminus}--\eqref{eq:global-Aplus}, the block-triangular dlog
potential \eqref{eq:global-diagonal-omega}--\eqref{eq:global-offdiagonal-omega}, and the
nested-blow-up boundary data, without generating or solving a large IBP system. The same registry
feeds the numerical backend FlintNDE, which transports the masters from the $K_v\to\infty$
boundary to user-supplied ordinary points.

\bibliographystyle{unsrtnat}
\bibliography{references}

\end{document}
